% =====================================================================
\chapter{Ángulos, ortogonalidad y similitud coseno}
\label{tema:2}
% =====================================================================

La desigualdad de Cauchy--Schwarz permite definir el ángulo entre dos vectores en cualquier dimensión. Este tema desarrolla esa geometría: ortogonalidad, proyección sobre una recta y similitud coseno, que es la medida con la que se comparan los embeddings. Se sigue la sección~1.2 de Strang~\cite{strang} y las secciones~3.3, 3.4 y~3.8 de~\cite{mml}; el comportamiento en alta dimensión procede del capítulo~2 de Blum, Hopcroft y Kannan~\cite{bhk}.

\section{El ángulo entre dos vectores}

\begin{definicion}[Ángulo]\label{def:angulo}
Sean $u,v\in\R^{d}$ no nulos. El \emph{ángulo} entre $u$ y $v$ es el único $\theta\in[0,\pi]$ tal que
\[
\cos\theta=\frac{u\cdot v}{\norm{u}\,\norm{v}}.
\]
\end{definicion}

La definición es correcta porque, por el \cref{teo:cs}, el cociente está en $[-1,1]$, y la función coseno es una biyección de $[0,\pi]$ en $[-1,1]$. La proposición siguiente muestra que, en el plano, coincide con el ángulo de la geometría elemental.

\begin{proposicion}\label{prop:ley-coseno}
Sean $u,v\in\R^{2}$ no nulos y sea $\varphi$ el ángulo geométrico que forman las flechas $u$ y $v$. Entonces $u\cdot v=\norm{u}\norm{v}\cos\varphi$.
\end{proposicion}

El signo del producto escalar indica, por tanto, si dos vectores apuntan en direcciones parecidas u opuestas (\cref{tab:signo}). El ejercicio~\ref{eje:2-tres} explora hasta dónde pueden «llevarse mal» varios vectores del plano.

\begin{table}[h]
\centering\small
\caption{Signo del producto escalar y ángulo entre dos vectores no nulos.}\label{tab:signo}
\begin{tabular}{@{}lll@{}}
\toprule
Producto escalar & Ángulo & Interpretación \\
\midrule
$u\cdot v>0$ & $0\le\theta<\pi/2$ & direcciones parecidas \\
$u\cdot v=0$ & $\theta=\pi/2$ & direcciones perpendiculares \\
$u\cdot v<0$ & $\pi/2<\theta\le\pi$ & direcciones opuestas \\
\bottomrule
\end{tabular}
\end{table}

\begin{definicion}[Ortogonalidad]
Dos vectores $u,v\in\R^{d}$ son \emph{ortogonales}, y se escribe $u\perp v$, si $u\cdot v=0$.
\end{definicion}

El vector nulo es ortogonal a todos, y es el único vector ortogonal a sí mismo, por la positividad del producto escalar. Los vectores de la base canónica son ortogonales dos a dos: $e_i\cdot e_j=0$ si $i\neq j$.

\begin{teorema}[Pitágoras]\label{teo:pitagoras}
Si $u\perp v$, entonces $\norm{u+v}^{2}=\norm{u}^{2}+\norm{v}^{2}$.
\end{teorema}

\section{Proyección sobre una recta}

\begin{proposicion}[Proyección ortogonal sobre una recta]\label{prop:proyeccion}
Sean $u,a\in\R^{d}$ con $a\neq 0$. Entre todos los vectores de la recta $\{ta: t\in\R\}$, el más próximo a $u$ es
\[
p=\frac{u\cdot a}{a\cdot a}\,a,
\]
y es el único punto de la recta tal que $u-p\perp a$. El número $(u\cdot a)/\norm{a}$ se llama \emph{proyección escalar} de $u$ sobre $a$; si $\norm{a}=1$, coincide con $u\cdot a$.
\end{proposicion}

\addtocounter{figure}{1}%

\addtocounter{definicion}{1}%

La proyección sobre una recta es el caso más sencillo de la proyección sobre un subespacio, que conduce a los mínimos cuadrados~\cite[§4.2]{strang}; en el \cref{tema:4} se escribe como producto de matrices (ejercicio~\ref{eje:4-proyeccion}).

\section{Similitud coseno}

\begin{definicion}[Similitud coseno]\label{def:coseno}
La \emph{similitud coseno} de dos vectores no nulos $u,v\in\R^{d}$ es
\[
\cos(u,v)=\frac{u\cdot v}{\norm{u}\,\norm{v}}=\hat u\cdot\hat v\in[-1,1].
\]
\end{definicion}

\begin{proposicion}\label{prop:coseno}
Sean $u,v\in\R^{d}$ no nulos.
\begin{enumerate}
  \item $\cos(au,bv)=\cos(u,v)$ para todos $a,b>0$, y $\cos(u,-v)=-\cos(u,v)$.
  \item $\cos(u,v)=1$ si y solo si $u=cv$ con $c>0$; $\cos(u,v)=-1$ si y solo si $u=cv$ con $c<0$.
\end{enumerate}
\end{proposicion}

La similitud coseno mide, por tanto, solo la dirección: $\cos((1,2),(2,4))=1$ aunque los vectores tengan longitudes distintas.

\addtocounter{definicion}{1}%

\begin{proposicion}[Coseno y distancia]\label{prop:coseno-distancia}
Si $\norm{u}=\norm{v}=1$, entonces $\norm{u-v}^{2}=2-2\cos(u,v)$. En consecuencia, entre vectores normalizados, ordenar los candidatos por distancia creciente a una consulta equivale a ordenarlos por similitud coseno decreciente.
\end{proposicion}

La demostración se propone en el ejercicio~\ref{eje:2-distancia}. Sin normalizar, ambas medidas pueden discrepar: un vector muy largo en la misma dirección que la consulta tiene coseno~1 pero puede estar lejos (ejercicio~\ref{eje:2-vecinos}). Los buscadores vectoriales de los sistemas RAG (módulo~22) usan el coseno o el producto escalar sobre vectores normalizados.

\section{Geometría en alta dimensión}

En $\R^{2}$, dos direcciones elegidas al azar forman ángulos de cualquier tamaño. En dimensión alta el comportamiento es muy distinto: dos direcciones aleatorias son casi ortogonales~\cite[§2.4.2]{bhk}. El resultado siguiente lo cuantifica.

\begin{proposicion}\label{prop:alta-dim}
Sea $U$ un vector aleatorio de $\R^{d}$ con $\norm{U}=1$ cuya distribución no cambia al permutar sus componentes ni al cambiar el signo de cualquiera de ellas (por ejemplo, $U=G/\norm{G}$ con $G_1,\dots,G_d$ normales estándar independientes). Para todo vector unitario fijo $v$,
\[
\E[U\cdot v]=0,\qquad \E\bigl[(U\cdot v)^{2}\bigr]=\frac1d .
\]
En particular, la similitud coseno entre $U$ y $v$ tiene media~0 y desviación típica $1/\sqrt d$.
\end{proposicion}

Si $v$ también es aleatorio e independiente de $U$, el resultado se aplica condicionando al valor de $v$. Para $d=768$, la desviación típica del coseno entre dos direcciones aleatorias es $1/\sqrt{768}\approx 0{,}0361$; para $d=12\,288$, $\approx 0{,}0090$. Una simulación con 2000 pares de vectores gaussianos da $0{,}0363$ y $0{,}0091$, respectivamente. Por eso un coseno de $0{,}3$ entre dos embeddings de dimensión 768, que equivale a más de ocho desviaciones típicas, es una señal fuerte de parecido.

\begin{observacion}
Este fenómeno permite que un espacio de dimensión $d$ contenga muchas más de $d$ direcciones casi ortogonales entre sí. La hipótesis de \emph{superposición} de la interpretabilidad (módulo~23) se apoya en esta propiedad: un modelo puede representar más conceptos que dimensiones tiene su espacio, asignando a cada uno una dirección casi ortogonal a las demás.
\end{observacion}

